
\subsubsection{3.1 Data Source and Model Set}

Published evaluation scores from \textbf{TrustLLM} \citep{Sun2024TrustLLM} (arXiv:2401.05561) are used throughout. TrustLLM evaluates 16 decoder-only LLMs using consistent scoring protocols on all target dimensions. Using a single source eliminates potential confounds from cross-protocol aggregation.

The 16 models span a broad capability range (MMLU: 0.278--0.864): GPT-4, GPT-3.5-Turbo, LLaMA-2-70B-Chat, LLaMA-2-13B-Chat, LLaMA-2-7B-Chat, LLaMA-2-70B, LLaMA-2-13B, LLaMA-2-7B, Mistral-7B-Instruct, Mistral-7B, Falcon-40B, Falcon-7B, Vicuna-13B, Alpaca-13B, Koala-13B, OpenAssistant-12B. Three models --- Alpaca-13B, Koala-13B, and OpenAssistant-12B --- lack robustness scores in TrustLLM; robustness analyses use N = 13.

\subsubsection{3.2 Benchmark Pairs and Distribution Shift Types}

\begin{table*}[htbp]
\centering
\resizebox{\dimexpr 2\columnwidth+\columnsep\relax}{!}{%
\begin{tabular}{lllll}
\toprule
Dimension & In-Distribution & Out-of-Distribution & Shift Type & N \\
\midrule
Fairness & BBQ-Disambig & BBQ-Ambig & Context informativeness & 16 \\
Robustness (Adversarial) & ANLI R1 & ANLI R3 & Adversarial difficulty & 13 \\
Robustness (OOD) & In-distribution NLU & OOD robustness (TrustLLM Micro F1) & Multi-condition OOD & 13 \\
\bottomrule
\end{tabular}
}
\end{table*}

The GLUE $\rightarrow$ AdvGLUE pair was not included: GLUE-X \citep{Yang2023GLUEX} evaluates encoder-only PLMs exclusively, with zero model overlap with the decoder-only LLMs in TrustLLM. The OOD robustness arm from TrustLLM serves as a proxy for the AdvGLUE analysis.

Score ranges across the 16-model set: BBQ-Disambig 0.33--0.89; BBQ-Ambig 0.45--0.72; ANLI R1 0.34--0.62; ANLI R3 0.33--0.57; OOD Robustness (Micro F1) 0.37--0.77; MMLU 0.28--0.86.

\subsubsection{3.3 Partial Spearman $\rho$ (MMLU-Controlled)}

\textbf{Motivation.} Raw Spearman $\rho$ between in-distribution and OOD model rankings is near-ceiling across all dimensions ($\rho \approx$ 0.978--0.984). This near-ceiling behavior reflects a strong general capability confound: models with higher general ability tend to score higher on all benchmarks. Partial correlation removing MMLU rank isolates dimension-specific trustworthiness structure.

\textbf{Computation.} Partial Spearman $\rho$ is computed via \texttt{pingouin.partial\textbackslash \_corr} (method='spearman', alternative='greater', $\alpha$ = 0.05). The fairness-robustness differential ($\Delta\rho$) is tested using the Fisher z-test for dependent correlations \citep{Meng1992}, applied to the N = 13 subset for both fairness and robustness to satisfy the same-sample assumption. For this test, $\rho_{fairness}$ is re-estimated on the N = 13 subset ($\rho$ = 0.967), yielding $\Delta\rho$ = 0.967 - 0.776 = 0.192. The N = 16 fairness estimate ($\rho$ = 0.962) is the primary result for RQ1.

\textbf{Sensitivity analysis.} The fairness analysis is repeated with Winogrande as the capability covariate (N = 14, two models lacking Winogrande scores excluded).

\subsubsection{3.4 Rank Reversal Analysis}

For each adversarial robustness pair, we count model pairs (A, B) where A outranks B on the in-distribution benchmark but B outranks A on the OOD benchmark. Zero reversals would indicate complete rank preservation; high counts would indicate disruption.

\subsubsection{3.5 Implementation}

Python 3.10.20; pingouin 0.6.1; scipy; pandas; matplotlib. All analyses validated through 9 unit tests (h-m1) and 6 mechanism indicator checks (h-m3). No model training or new data collection was required.

